\section{深入认识前端技术}
\subsection{梳理网页组成}
开始这次任务时，我对前端只有一些粗浅的了解，比如说：知道 three.js 能够帮助网站实现三维效果，但却不知道具体如何使用等等。真正接触项目后，我才第一次从代码层面接触了首页布局、%
菜单展开、文章排版、图片展示和后台编辑等多种多样的前端模块。%
因此，我首先把前端技术按解决的问题分类，然后再决定应该使用什么样的前端工具来完成本项目。

通过调研我了解到：HTML 负责表达内容结构，CSS 负责布局与视觉样式，JavaScript 负责交互与逻辑，%
这是我建立的第一层认识\cite{mdn-basics}。以一张文章卡片为例，标题、图片和链接是内容结构；%
卡片宽度、字距与颜色是样式；点击筛选后重新获取文章是交互。这三者可以类比为展览中的展品目录、%
空间布置与现场工作人员，它们在浏览器中共同作用，共同呈现出一个完整的可交互界面。

浏览器把 HTML 解析成 DOM，即文档对象模型，可以将其理解为一棵由标题、段落、图片等节点组成的树。%
JavaScript 通过 DOM API 读取或修改这些节点；CSS 则影响它们的尺寸、位置和显示方式。%
理解这一点后，我才逐渐明白“把文字放到网页上”和“让网页持续回应用户操作”之间的巨大区别：前者只需呈现内容，%
后者还需要管理事件、数据与页面状态。

\begin{figure}[H]
\centering
\begin{tikzpicture}[font=\fontsize{11}{15}\selectfont,
 box/.style={draw=wmaccent,fill=wmnavy!4,rounded corners=2pt,%
   text width=13.4cm,minimum height=1.15cm,align=center},
 arr/.style={-{Stealth[length=2mm]},draw=wmaccent,line width=0.7pt}]
\node[box] (a) at (0,0) {\textbf{基础：网页有什么，怎样呈现，如何交互}\\%
  HTML 结构 \quad CSS 样式 \quad JavaScript 行为};
\node[box] (b) at (0,-1.6) {\textbf{组织：页面多了以后如何协作与复用}\\%
  Vue 组件 \quad Router 路由 \quad Pinia 共享状态};
\node[box] (c) at (0,-3.2) {\textbf{工程：怎样开发、检查并交付}\\%
  TypeScript 类型 \quad Vite 构建 \quad 测试与浏览器检查};
\node[box,fill=wmaccent!12] (d) at (0,-4.8) {\textbf{任务扩展：如何改善展示和内容操作}\\%
  三维与动效 \quad 内容与图片加载 \quad 图文编辑};
\draw[arr] (a)--(b);%
\draw[arr] (b)--(c);%
\draw[arr] (c)--(d);
\end{tikzpicture}

\caption{我在任务中逐步建立的前端技术地图}%
\label{fig:map}
\end{figure}

\subsection{从基础语言到框架和工程工具}
第二层认识是：框架并不取代 HTML、CSS 和 JavaScript，而是帮助组织它们。%
Vue 的单文件组件将模板、逻辑与样式放在一个组件中，配合响应式状态驱动界面更新\cite{vue-intro}。%
例如，文章列表组件接收数据并显示卡片，公共导航组件统一处理菜单；我不必在每个页面重新写一套相同结构。

我也调研了 React、Svelte 等方案。React 同样以组件组织界面，强调组件、状态和数据流；%
Svelte 则利用编译器转换组件代码\cite{react-ui,svelte-overview}。这些框架的共同点是减少手工维护复杂
DOM 的负担，但语法、状态机制和配套生态不同。%


TypeScript 和 Vite 又属于不同层次。TypeScript 为 JavaScript 增加静态类型检查，%
便于发现字段或参数使用错误；类型标注在编译后会被移除，不能自动验证服务器实际返回的
JSON\cite{ts-basics}。Vite 提供开发服务器、模块更新和生产构建流程，%
解决的是开发与交付效率\cite{vite-why}。

\subsection{基础技术分类}
\begin{table}[!htb]
\centering
\smalltable
\caption{前端技术分别解决什么问题}
\begin{tabularx}{\linewidth}{>{\raggedright\arraybackslash}p{2.4cm}YY}
\toprule
层次 & 代表技术 & 我在项目中对应的理解\\
\midrule
内容与样式 & HTML、CSS、Flexbox、Grid & 组织文章、排列导航和卡片，适应不同页面宽度\\
\tabledivider
行为与平台能力 &
    JavaScript、DOM、Fetch、事件、动画 API &
    读取数据、响应点击、控制加载状态和动态效果\\
\tabledivider
组件与状态组织 & Vue、Router、Pinia & 复用页面结构，切换地址，共享必要的会话状态\\
\tabledivider
工程与质量 & TypeScript、Vite、Vitest & 提前发现错误、构建资源、检查关键交互逻辑\\
\tabledivider
专题展示与内容生产 & Three.js、图片加载、富文本编辑器 & 为具体体验补充能力，而非每个页面全部引入\\
\bottomrule
\end{tabularx}
\end{table}

这次分类还帮助我区分了两个容易混淆的“响应式”：Vue 的响应式是数据变化后界面跟着更新，CSS
的响应式设计是布局随可用空间变化。前者主要处理数据与视图的关系，后者处理内容与屏幕的关系。%

